\section{Related Work}
\label{sec:related}

Our work sits at the intersection of multi-dimensional LLM trustworthiness evaluation, adversarial robustness benchmarking, and cross-benchmark predictive validity methodology. We describe each line and explain what it leaves unanswered.

\subsection{Multi-Dimensional LLM Trustworthiness Evaluation}

\textbf{TrustLLM} \citep{Sun2024TrustLLM} is the most comprehensive publicly available trustworthiness evaluation of open and proprietary LLMs, covering 16 models across six dimensions including fairness (BBQ), robustness (ANLI, OOD tasks), privacy, and safety. TrustLLM provides consistent single-source evaluation scores that enable our analysis, but does not compute rank correlation between in-distribution and OOD benchmark variants as a research question. It reports absolute scores, not predictive validity.

\textbf{DecodingTrust} \citep{Wang2023DecodingTrust} evaluates GPT-3.5 and GPT-4 on eight trustworthiness dimensions including adversarial robustness and fairness, finding that GPT-4 is more vulnerable to adversarial jailbreaking despite higher standard benchmark scores. This $N=2$ rank reversal observation motivated our adversarial disruption hypothesis --- which our $N=16$ analysis falsifies. The lesson is that two-model anecdotes do not scale.

\textbf{HELM} \citep{Liang2022Holistic} provides holistic multi-scenario evaluation across 42 models and 7 metrics, establishing the multi-benchmark multi-model evaluation paradigm. Like TrustLLM and DecodingTrust, HELM reports within-scenario scores rather than asking whether one scenario's rankings predict another's. Our predictive validity framing is orthogonal to --- and enabled by --- this evaluation infrastructure.

These evaluations collectively establish that trustworthiness is multi-dimensional and model-specific. What they do not provide is an analysis of whether evaluation results on in-distribution conditions transfer predictively to OOD conditions.

\subsection{Adversarial Robustness Benchmarks}

\textbf{AdvGLUE} \citep{Wang2021AdvGLUE} applies 14 adversarial attack methods to GLUE tasks, producing a benchmark explicitly designed to defeat models that pass the in-distribution GLUE tasks. The design philosophy --- iterative human-model adversarial data collection to maximally challenge current models --- underpinned our initial hypothesis that rank stability would be disrupted. Our results show that OOD robustness rank stability ($\rho = 0.868$) contradicts this intuition at the model-ranking level.

\textbf{ANLI} \citep{Nie2020Adversarial} constructs adversarial NLI rounds (R1, R2, R3) through iterative human-and-model-in-the-loop processes, with each round designed to defeat models that pass earlier rounds. Our finding that partial $\rho_{\text{ANLI}} = 0.684$ ($p = 0.007$) after MMLU control is positive indicates that adversarial construction disrupts performance but preserves relative model ordering.

\textbf{Wang et al.}\ \citep{Wang2023Robustness} evaluate ChatGPT on AdvGLUE and ANLI against baseline models, finding ChatGPT advantages on OOD tasks but not computing cross-model rank correlations or controlling for capability. Our study extends this to 16 models with capability control.

\subsection{Fairness Benchmark Design}

\textbf{BBQ} \citep{Parrish2022BBQ} constructs 50,000 QA items across 9 social bias categories, with disambiguated contexts (where factual information resolves the answer) and ambiguous contexts (where only stereotype-consistent or stereotype-inconsistent guessing is possible). The disambiguated$\rightarrow$ambiguous shift constitutes a natural in-distribution/OOD pair for fairness. Our finding that they share near-perfect partial $\rho = 0.962$ validates BBQ's construct validity at the cross-model level while raising the question of whether this reflects stable model mechanisms or shared benchmark structure.

\subsection{Cross-Benchmark Predictive Validity}

\textbf{Gevers and Daelemans} \citep{GeversAndDaelemans2026} provide the closest methodological predecessor: rank correlations with leave-one-family-out cross-validation across 23 LLMs on commonsense benchmarks (preprint; independently verified as plausible but not confirmed via library access at submission time). We extend their approach to trustworthiness dimensions, which differ from commonsense tasks in (1) higher safety stakes motivating the analysis, and (2) an ID/OOD structure reflecting mechanistically distinct properties (stable latent biases for fairness vs.\ adversarial design for robustness).

\textbf{GLUE-X} \citep{Yang2023GLUEX} measures in-distribution to OOD accuracy gaps for encoder-only PLMs, finding systematic degradation. GLUE-X does not compute cross-model rank correlations and covers encoder-only PLMs --- not the decoder-only LLMs in our study. The architectural divergence is why GLUE$\rightarrow$AdvGLUE was unavailable for our LLM analysis.

No prior study computes capability-controlled partial Spearman $\rho$ between in-distribution and OOD trustworthiness benchmark model rankings across multiple dimensions and 15+ models. We fill this gap.
